\documentclass[10pt]{extarticle}

% ---------- Packages ----------
\usepackage{fancyhdr}
\usepackage{geometry}
\usepackage{hyperref}
\usepackage{xcolor}
\usepackage{amsmath, amsfonts, bm}
\usepackage{subcaption}
\usepackage{graphicx}
\usepackage{tikz}
\usepackage{nomencl}
\usepackage{array}
\usepackage{pdflscape}
\usepackage[export]{adjustbox}
\usepackage[utf8]{inputenc}
\usepackage{parskip}
\usepackage{titlesec}

\usetikzlibrary{shapes.geometric}
\usetikzlibrary{arrows.meta, arrows}

% ---------- Page setup ----------
\geometry{
    a4paper,
    total={170mm,257mm},
    left=20mm,
    top=20mm
}

% ---------- Headers and footers ----------
\pagestyle{fancy}
\pagenumbering{arabic}

\fancyhead[L]{Bike Physics by Yaroslav Knyazkov}
\fancyfoot[C]{}
\fancyfoot[R]{\thepage}

% ---------- Hyperlinks ----------
\hypersetup{
    colorlinks=true,
    linkcolor={blue!50!black},
    citecolor={blue!50!black},
    urlcolor={blue!80!black}
}

% ---------- Equation numbering ----------
\numberwithin{equation}{subsection}

% ---------- Custom commands ----------
\newcommand{\todo}[1]{%
    \textcolor{red}{(#1)}%
}
% \ai{...} marks short inline text written/edited by Claude (blue). Do not put blank lines
% or sectioning commands inside it -- \textcolor's argument can't contain a \par.
\newcommand{\ai}[1]{\textcolor{blue}{#1}}
% For multi-paragraph blocks (incl. \subsection etc.), use the aiblock environment instead:
% \begin{aiblock} ... \end{aiblock}
\newenvironment{aiblock}{\color{blue}}{}
% \vv{X} = the vector X, always rendered with a full-width arrow (\overrightarrow), for both
% single-symbol quantities (\vv{F}, \vv{x}) and named offset vectors (\vv{AB}, \vv{A_1A_{2a}}).
% Not called \v -- that's LaTeX's built-in caron accent (\v{c} produces c-with-caron).
\newcommand{\vv}[1]{\overrightarrow{#1}}
\newcommand{\sgn}{\operatorname{sgn}}

% ---------- 4th heading level: System -> Subsystem instance -> content category ----------
% \subsection   = System        (Frame + Rider / Wheel / Drivetrain)
% \subsubsection = Subsystem instance (Rear (2a), Front (2b), Cassette (3a), Crankset (3b));
%                  omitted for single-instance systems like Frame + Rider.
% \subsubsubsection = content category (Assumptions / Force / Geometry / Moment)
% Unnumbered by design: these labels repeat identically under every instance, a running
% number would add noise without adding information.
\titleformat{\paragraph}[block]{\normalfont\bfseries}{}{0pt}{}
\titlespacing*{\paragraph}{0pt}{1.5ex plus 1ex minus .2ex}{0.5ex plus .1ex}
\newcommand{\subsubsubsection}[1]{\paragraph{#1}}

\begin{document}

\tableofcontents
\newpage

\section{Timeline}
This is a roadmap for the project as a whole, not just this document. Four future application tracks are planned, on top of the physics core work already in progress here.

\begin{enumerate}
    \item \textbf{Physics simulation} (core, ongoing; currently the main focus). A numerical implementation of the multibody model derived in this document, itself built in three tiers -- see \S\ref{sec:simplified}--\S\ref{sec:3d} below for what each contains:
    \begin{itemize}
        \item \emph{Simplified}: closed-form, analytically-tractable models (forward motion, turning motion). Supports and sanity-checks the numeric stages below rather than being superseded by them.
        \item \emph{2D}: the app's first real, coded-system-by-system stage -- current focus. Build order:
        \begin{enumerate}
            \item (in progress) simple frame, wheels, simplified (lumped) drivetrain
            \item complex drivetrain -- each component (chain, cassette, crankset, pedals, \dots) as its own body
            \item suspension -- front fork and rear, with multiple rear-suspension kinematic configurations
            \item tire stiffness -- relax the rigid-wheel assumption
            \item rider legs isolated -- full leg model with actuators, still attached to the frame
            \item rider fully isolated -- whole rider as its own body, including postural/jump dynamics
        \end{enumerate}
        Other candidate refinements, not yet prioritized against the list above: see Appendix~\ref{app:refinements}.
        \item \emph{3D}: final stage, much later. Initially motion only, extending the 2D multibody model to full spatial dynamics; other purposes may layer on afterward. Priority is to finish 2D forward-motion and all four application tracks on top of it first.
    \end{itemize}
    \item \textbf{Optimization and strategy} (primary goal). Two related but distinct problems built on top of the simulation core:
    \begin{itemize}
        \item \emph{Optimization}: given an objective (performance / time / velocity) and constraints, find the input (rider effort profile, line, gearing, \dots) that achieves it. Target order: road bike first (most predictable/constrained), then BMX, then MTB (least predictable -- terrain and technical variability dominate).
        \item \emph{Strategy}: combine an energy model (rider/team energy expenditure and availability) with a strategy model (tactical choices -- drafting, attacks, team formations) to generate multiple possible race outcomes, the way F1 strategy simulation does -- surfacing which choices let a rider or team win, or maximize some other goal, rather than solving a single fixed objective directly. Planned starting point: define a rider power curve (or equivalent) as the energy-model primitive.
    \end{itemize}
    \item \textbf{Research and development}. A tool for bike manufacturers: given a desired behavior, or a trade-off to explore (e.g. weight vs.\ aero), use the model to steer component/geometry design. Two concrete deliverables planned: (a) a set of open questions the model should be able to answer analytically or numerically, and (b) a systematic map from bike geometry parameters to the behavior each one drives.
    \item \textbf{Data acquisition}. Let ordinary riders, without lab equipment (wind tunnel, dynamometer, etc.), estimate their own real-world parameters -- drag ($C_dA$), rolling/bearing friction, drivetrain efficiency, braking coefficient -- from ordinary ride data. This is also how the simulation core in (1) gets calibrated against reality, so it feeds every other track rather than standing apart from them.
\end{enumerate}

Rough dependency: (1) underlies (2) and (3); (4) calibrates (1) and is largely independent otherwise, so it can be worked on in parallel.


\section{Description}
This document holds the analytical derivations backing the app's code -- not a standalone reference, but the math worked out before it gets implemented. It has three parts, matching the Physics-simulation tiers in the Timeline above:
\begin{itemize}
    \item \textbf{Simplified}: closed-form models, simplified enough to stay analytically tractable (forward motion, turning motion). Not started yet.
    \item \textbf{2D}: full planar (sagittal-plane) rigid-body multibody derivation, 3 DOF per body (see \texttt{Current global assumptions} below) -- current focus, and the only part with real content so far.
    \item \textbf{3D}: full spatial rigid-body derivation, 6 DOF per body. Future work.
\end{itemize}

The 2D part is organized bottom-up:
\begin{itemize}
    \item \textbf{2D -- Systems}: free-body force/moment equations for each rigid body in isolation (Frame + Rider, Wheel -- Rear/Front, Drivetrain -- Cassette/Crankset), following a fixed template per subsystem instance: Assumptions, Force, Geometry, Moment.
    \item \textbf{2D -- Assembly}: how the isolated systems connect -- joint reaction pairs and kinematic constraints (Frame -- Wheels), and the ground-contact model (Wheels -- Ground), including the curved-terrain and stick/slip/airborne cases needed for the full range of bike dynamics.
\end{itemize}
Conventions: subscripts follow the numbering in \texttt{Systems} (1 = Frame + Rider, 2a/2b = Rear/Front wheel, 3a/3b = Cassette/Crankset); \todo{...} marks open gaps in red; \texttt{aiblock} (or inline \texttt{\textbackslash ai\{...\}}) marks text written or edited by Claude in blue, so it stays distinguishable from your own writing.


\newpage
\section{Simplified}
\label{sec:simplified}
\begin{aiblock}
Closed-form, analytically-tractable models -- simplified enough to solve by hand or in closed form. Support and sanity-check the numeric 2D/3D work rather than being superseded by it.
\end{aiblock}

\subsection{Forward motion}
\label{sec:forward-motion}
\begin{aiblock}
\ai{Ported from a standalone draft (\texttt{temp/simplified.tex}, since removed) and translated into this project's conventions: body indices $1$ (Frame+Rider), $2a$/$2b$ (Rear/Front wheel), $3b$ (Crankset, cassette folded into $2a$ per \texttt{wheel.tex}); frame $i$ (fixed) / $0$ (ground-local); $\theta$ for ground/frame tilt, $\phi$ for body spin (was $\theta_1$ and $\theta_{2,3,4}$ respectively -- the original's road-gradient $\theta_1$ collided with this project's spin-angle symbol once translated, hence the split). Point labels ($A$--$K$, $a$) are local to this derivation and are \emph{not} the same points as in 2D -- Systems (e.g. $C$ here is the rear wheel's shifted ground contact, not the drivetrain bearing point in \texttt{frame.tex}). Air density $\rho$ in \S Analysis is unrelated to the road curvature radius $\rho(x)$ in 2D -- Assembly; both are standard symbols in their own contexts. $M$ is the total lumped system mass (frame+rider+wheels+drivetrain combined, per the point-mass assumption below) -- distinct from a per-body $m_k$. Equation numbers restart at each of Bike, Wheels, Drivetrain, Assembly, Analysis, per the document's \texttt{numberwithin\{equation\}\{subsection\}}.}

\ai{\textbf{Axis-direction note:} within this derivation only, $\hat z_i$ and $\hat z_0$ point \emph{downward} (along gravity), the opposite of 2D's global-assumption convention (where $\hat z_i$ points up, opposite gravity). This is inherited unchanged from the original draft and kept as-is rather than resigning the whole derivation, since every equation below is internally consistent under it -- e.g.\ $\vv W_1=+Mg\hat z_i$ and the Limits condition ${}^{z_0}R_{2a/0}\le0$ ("no wheel lift") only both hold together if $z$ points down. If this gets confusing in practice, the fix is a single global sign flip on every $z$-component here, not a symbol rename.}
\end{aiblock}

\begin{aiblock}
\subsection{Bike}
The equation of motion is written for the total bike system. Internal forces between subsystems cancel and do not appear.

\subsubsubsection{Assumptions}
\begin{itemize}
    \item Rigid frame --- no compliance, no rider motion relative to frame
    \item Point mass at fixed height $h$, total mass $M$ --- rider and frame CG collocated and invariant
    \item Motion confined to the $x$--$z$ plane --- straight-line travel, no yaw or roll
    \item Constant road gradient $\theta$ --- no curvature or undulation
    \item Two contact points only --- front and rear tyre, no additional supports
    \item Aerodynamic lift neglected --- drag and pitching moment only
    \item Wind in the $x$--$z$ plane --- no crosswind
    \item Rolling resistance as contact resultant shift (Option A) --- no separate resistance moment
    \item Both tyres share the same $c_{2a,rr}=c_{2b,rr}$ and $r_{2a}=r_{2b}$
\end{itemize}

\subsubsubsection{Forces}
Weight acts at the system centre of gravity $a$ in the fixed frame:
\begin{equation}
    \vv{W}_1 = Mg\,\hat{z}_i
\end{equation}

Under load, tire deformation shifts the contact resultant forward of the geometric contact by $l_{2a,rr} = c_{2a,rr}r_{2a}$ (rear), $l_{2b,rr} = c_{2b,rr}r_{2b}$ (front). Points $C$ and $D$ denote these shifted contact locations. The ground reactions at $C$ and $D$ are decomposed in the ground-local frame $(\hat{x}_0, \hat{z}_0)$:
\begin{gather}
    \vv{R}_{2a/0} = {}^{x_0}R_{2a/0}\,\hat{x}_0 + {}^{z_0}R_{2a/0}\,\hat{z}_0 \\
    \vv{R}_{2b/0} = {}^{x_0}R_{2b/0}\,\hat{x}_0 + {}^{z_0}R_{2b/0}\,\hat{z}_0
\end{gather}

Aerodynamic drag and pitching moment act on the frame, neglecting lift:
\begin{gather}
    \vv{A}_{1,f} = -A_{1,x}\,\hat{x}_0 \\
    \vv{A}_{1,m} = -A_{1,m}\,\hat{y}
\end{gather}

\subsubsubsection{Geometry}
Horizontal distances $l_C$, $l_D$ are measured from the system CG $a$ to the geometric wheel contact points. The rolling resistance offsets $l_{2a,rr}$, $l_{2b,rr}$ shift the contact resultant forward; $\sigma$ ensures this shift is always in the direction of motion:
\begin{gather}
    L = l_C + l_D \qquad \sigma = \sgn(\dot{x}^{(1)}) \\
    l_C' = l_C - l_{2a,rr}\sigma \qquad l_D' = l_D + l_{2b,rr}\sigma
\end{gather}
The position vectors from $a$ to each force application point, expressed in the ground-local frame:
\begin{gather}
    \vv{aC} = (-l_C + l_{2a,rr}\sigma)\,\hat{x}_0 + h\,\hat{z}_0 \\
    \vv{aD} = (+l_D + l_{2b,rr}\sigma)\,\hat{x}_0 + h\,\hat{z}_0 \\
    \vv{aA} = -l_C\,\hat{x}_0 + (h - r_{2a})\,\hat{z}_0
\end{gather}

\subsubsubsection{Equations of motion}
The moment equation uses $\dot{\vv{H}}_a$, the rate of change of angular momentum about $a$, rather than a single body inertia, because the spinning wheels and crankset each contribute angular momentum independently:
\begin{align}
    M\,\ddot{\vv{x}}^{(1)} &= \sum \vv{F}_1 = \vv{W}_1 + \vv{R}_{2a/0} + \vv{R}_{2b/0} + \vv{A}_{1,f} \\[6pt]
    \dot{\vv{H}}_a &= \sum \vv{M}_1(a) = \vv{aC}\times\vv{R}_{2a/0} + \vv{aD}\times\vv{R}_{2b/0} + \vv{aA}\times\vv{A}_{1,f} + \vv{A}_{1,m}
\end{align}

\subsubsubsection{Translation}
Weight is projected onto the ground-local frame using the gradient angle $\theta$. The vertical equation enforces the constraint that the bike remains on the road:
\begin{align}
    M\,\ddot{x}^{(1)} &= -Mg\sin\theta + {}^{x_0}R_{2a/0} + {}^{x_0}R_{2b/0} - A_{1,x} \\
    0                 &= Mg\cos\theta + {}^{z_0}R_{2a/0} + {}^{z_0}R_{2b/0}
\end{align}

\subsubsubsection{Rotation}
Since all subsystems translate together, the transport terms in $\dot{H}_a$ vanish and only the spin inertias remain. The frame's own rotation is zero ($\ddot\theta = 0$, since $\theta$ is held constant here), so:
\begin{align}
    \dot{H}_y &= I_{2a}\ddot{\phi}_{2a} + I_{2b}\ddot{\phi}_{2b} + I_{3b}\ddot{\phi}_{3b} \\
    \dot{H}_y &= h\left({}^{x_0}R_{2a/0} + {}^{x_0}R_{2b/0}\right) + l_C'\,{}^{z_0}R_{2a/0} - l_D'\,{}^{z_0}R_{2b/0}
               - \left(h - r_{2a}\right)A_{1,x} - A_{1,m}
\end{align}

\subsection{Wheels}
Each wheel is an independent rigid body. Moments are taken about the axle to eliminate the axle reaction, which passes through it and produces no moment.

\subsubsubsection{Assumptions}
\begin{itemize}
    \item Each wheel is a rigid body with uniform mass distribution --- CG at the axle
    \item Axle fixed relative to the frame --- no suspension travel
    \item Wheel spin axis is $\hat{y}$ --- no camber or toe
    \item Contact patch is a point --- no contact patch width or length effects
    \item No lateral forces --- pure longitudinal contact resultant in the $x$--$z$ plane
    \item Brake force acts at a single effective radius ($r_{2a,bk}$ rear, $r_{2b,bk}$ front) --- caliper geometry collapsed
    \item No wheel flex or tyre sidewall compliance
\end{itemize}

\subsubsubsection{Geometry}
Both wheels share the same radius ($r_{2a}=r_{2b}$) and rolling resistance coefficient ($c_{2a,rr}=c_{2b,rr}$), so the axle-to-contact vector has the same form for both. The chain and brake moment arms are specific to each wheel:
\begin{equation}
    |\vv{AF}| = r_{2a,cg} \qquad |\vv{AJ}| = r_{2a,bk} \qquad |\vv{BK}| = r_{2b,bk}
\end{equation}
The axle-to-contact resultant vector, expressed in the ground-local frame:
\begin{gather}
    \vv{AC} = l_{2a,rr}\sigma\,\hat{x}_0 + r_{2a}\,\hat{z}_0 \\
    \vv{BD} = l_{2b,rr}\sigma\,\hat{x}_0 + r_{2b}\,\hat{z}_0
\end{gather}

\subsubsubsection{Forces}
The chain tension acts only on the rear wheel at the cog contact point $F$. Brake forces act tangentially at the effective pad radius on whichever wheel is braked:
\begin{gather}
    \vv{T}_{2a/3b} = T_{2a/3b}\,\hat{e}_{span} \\
    \vv{B}_{2a/1} = B_{2a/1}\,\hat{e}_{t,bk}, \qquad \vv{B}_{2b/1} = B_{2b/1}\,\hat{e}_{t,bk}
\end{gather}

\subsubsubsection{Equations of motion}
Expanding the cross products with the geometry above, noting that $\vv{R}_{2a/1}$ and $\vv{W}_{2a}$ both act through the axle and vanish (likewise $\vv{R}_{2b/1}$, $\vv{W}_{2b}$):
\begin{align}
    I_{2a}\ddot{\phi}_{2a} &= \sum \vv{M}_{2a}(A)\cdot\hat{y}
        = \bigl[\vv{AC}\times\vv{R}_{2a/0} + \vv{AF}\times\vv{T}_{2a/3b}
          + \vv{AJ}\times\vv{B}_{2a/1}\bigr]\cdot\hat{y} \\[6pt]
    I_{2b}\ddot{\phi}_{2b} &= \sum \vv{M}_{2b}(B)\cdot\hat{y}
        = \bigl[\vv{BD}\times\vv{R}_{2b/0}
          + \vv{BK}\times\vv{B}_{2b/1}\bigr]\cdot\hat{y}
\end{align}

\subsubsubsection{Rear wheel}
The traction force ${}^{x_0}R_{2a/0}$ drives the wheel; the normal load ${}^{z_0}R_{2a/0}$ contributes a rolling resistance moment through the offset $l_{2a,rr}\sigma$; chain tension resists spin; braking resists spin:
\begin{equation}
    I_{2a}\ddot{\phi}_{2a} = r_{2a}\,{}^{x_0}R_{2a/0} - l_{2a,rr}\sigma\,{}^{z_0}R_{2a/0}
        - r_{2a,cg}\,T_{2a/3b} + \sigma\,r_{2a,bk}\,B_{2a/1}
\end{equation}

\subsubsubsection{Front wheel}
The front wheel carries no chain force. Only the contact resultant and the brake act:
\begin{equation}
    I_{2b}\ddot{\phi}_{2b} = r_{2b}\,{}^{x_0}R_{2b/0} - l_{2b,rr}\sigma\,{}^{z_0}R_{2b/0}
        + \sigma\,r_{2b,bk}\,B_{2b/1}
\end{equation}

\subsection{Drivetrain}
The drivetrain subsystem comprises the crankset: chainring, cranks, and pedals. The cassette is treated as part of the rear wheel (matches Rear (2a) already carrying the cassette reaction, per \texttt{wheel.tex}). A fixed gear ratio $r_{2a,cg}/r_{3b,cr}$ is assumed throughout.

\subsubsubsection{Assumptions}
\begin{itemize}
    \item Fixed gear ratio --- no derailleur shifting, single chainring and cog
    \item No drivetrain efficiency loss --- chain transmits torque without friction
    \item Chain mass neglected --- chain inertia not modelled
    \item Crankset is a rigid body with CG at the BB axle $E$
    \item Pedal force acts tangentially at the crank tip --- no radial component
    \item Left and right cranks combined as a single effective torque
    \item No freewheel modelled explicitly --- $T_{2a/3b} \ge 0$ enforced as a post-check
\end{itemize}

\subsubsubsection{Geometry}
$E$ is the bottom bracket axle, the rotation centre of the crankset. The chainring engages the chain at $G$; the rider applies force at the pedal spindles $H_L$, $H_R$. The relevant moment arms are:
\begin{equation}
    |\vv{EG}| = r_{3b,cr} \qquad |\vv{EH_L}| = |\vv{EH_R}| = r_{1,pedal}
\end{equation}

\subsubsubsection{Forces}
The chain pulls the chainring at $G$ with the same tension it pulls the cog, by Newton's third law. The rider applies pedal forces tangentially at each crank tip:
\begin{gather}
    \vv{T}_{3b/2a} = -\vv{T}_{2a/3b} \\
    \vv{P}_{3b/1,L} = P_{3b/1,L}\,\hat{e}_{t,L}, \qquad \vv{P}_{3b/1,R} = P_{3b/1,R}\,\hat{e}_{t,R}
\end{gather}

\subsubsubsection{Equations of motion}
Moments about $E$ eliminate the BB axle reaction and crankset weight:
\begin{equation}
    I_{3b}\ddot{\phi}_{3b} = \sum \vv{M}_{3b}(E)\cdot\hat{y}
        = \bigl[\vv{EG}\times\vv{T}_{3b/2a}
          + \vv{EH_L}\times\vv{P}_{3b/1,L}
          + \vv{EH_R}\times\vv{P}_{3b/1,R}\bigr]\cdot\hat{y}
\end{equation}

\subsubsubsection{Crankset}
Pedal force drives the crankset; chain tension resists it:
\begin{equation}
    I_{3b}\ddot{\phi}_{3b} = r_{3b,cr}\,T_{2a/3b} - r_{1,pedal}\left(P_{3b/1,L} + P_{3b/1,R}\right)
\end{equation}

\subsection{Assembly}
The four subsystems are coupled through three kinematic constraints: no-slip at each contact and the fixed chain ratio. All angular accelerations reduce to a single unknown $\ddot{x}^{(1)}$. The solve order has no algebraic loop: the EOM gives $\ddot{x}^{(1)}$ first, then reactions are post-processed in sequence.

\subsubsubsection{Assumptions}
\begin{itemize}
    \item No-slip at both contacts --- wheels roll without sliding
    \item Kinematic constraint is holonomic --- $\dot{x}^{(1)} = -r_{2a}\dot{\phi}_{2a} = -r_{2b}\dot{\phi}_{2b}$ at all times
    \item Chain does not slip on sprockets --- gear ratio is exact
    \item All assumptions from Bike, Wheels, and Drivetrain carry forward
\end{itemize}

\subsubsubsection{Kinematics}
No-slip at each contact relates wheel angular acceleration to frame linear acceleration. The chain ratio propagates this to the crankset:
\begin{align}
    \ddot{\theta} &= 0 \\
    \ddot{\phi}_{2a} = \ddot{\phi}_{2b} &= -\frac{\ddot{x}^{(1)}}{r_{2a}} \\
    \ddot{\phi}_{3b} &= -\frac{r_{2a,cg}}{r_{3b,cr}\,r_{2a}}\,\ddot{x}^{(1)}
\end{align}

\subsubsubsection{Chain tension}
Substituting kinematics into the crankset equation and solving for $T_{2a/3b}$. The first term captures crankset inertia; the second is the pedal input referred to the chain:
\begin{equation}
    \boxed{T_{2a/3b} = -\frac{I_{3b}\,r_{2a,cg}}{r_{3b,cr}^2\,r_{2a}}\,\ddot{x}^{(1)}
           + \frac{r_{1,pedal}}{r_{3b,cr}}\left(P_{3b/1,L} + P_{3b/1,R}\right)}
\end{equation}

\subsubsubsection{Traction}
Substituting kinematics into each wheel equation and solving for the longitudinal ground reaction. These are the horizontal forces the ground must supply; they are checked against the friction limit below:
\begin{align}
    \boxed{{}^{x_0}R_{2a/0} = -\frac{I_{2a}}{r_{2a}^2}\,\ddot{x}^{(1)}
        + c_{2a,rr}\sigma\,{}^{z_0}R_{2a/0}
        + \frac{r_{2a,cg}}{r_{2a}}\,T_{2a/3b}
        - \frac{\sigma\,r_{2a,bk}}{r_{2a}}\,B_{2a/1}} \\
    \boxed{{}^{x_0}R_{2b/0} = -\frac{I_{2b}}{r_{2b}^2}\,\ddot{x}^{(1)}
        + c_{2b,rr}\sigma\,{}^{z_0}R_{2b/0}
        - \frac{\sigma\,r_{2b,bk}}{r_{2b}}\,B_{2b/1}}
\end{align}

\subsubsubsection{Normal reactions}
The rotation equation, after substituting kinematics and traction, yields the auxiliary quantity $K$ which collects all contributions to the pitching moment. The normal reactions follow from $K$ and the vertical equilibrium. Note that the EOM needs only the sum of the normal loads, so this step is post-processing:
\begin{align}
    \boxed{{}^{z_0}R_{2a/0} = \frac{K - l_D'\,Mg\cos\theta}{L}} \qquad
    \boxed{{}^{z_0}R_{2b/0} = -\frac{K + l_C'\,Mg\cos\theta}{L}}
\end{align}
\begin{equation}
    K = -\frac{J}{r_{2a}}\,\ddot{x}^{(1)}
        - hM\!\left(\ddot{x}^{(1)} + g\sin\theta\right)
        - r_{2a}\,A_{1,x} + A_{1,m},
    \qquad
    J = I_{2a} + I_{2b} + \frac{r_{2a,cg}}{r_{3b,cr}}\,I_{3b}
\end{equation}

\subsubsubsection{Equation of motion}
Substituting traction into the bike translation equation and using ${}^{z_0}R_{2a/0} + {}^{z_0}R_{2b/0} = -Mg\cos\theta$ collapses the rolling resistance terms. The effective mass $M_{eq}$ reflects all rotational inertias to the frame linear acceleration:
\begin{equation}
    \resizebox{\linewidth}{!}{$\displaystyle
    \boxed{
    M_{eq}\,\ddot{x}^{(1)} =
        \frac{r_{1,pedal}\,r_{2a,cg}}{r_{3b,cr}\,r_{2a}}\left(P_{3b/1,L} + P_{3b/1,R}\right)
        - \frac{\sigma}{r_{2a}}\left(r_{2a,bk}\,B_{2a/1} + r_{2b,bk}\,B_{2b/1}\right)
        - Mg\sin\theta
        - \sigma\,c_{2a,rr}\,Mg\cos\theta
        - A_{1,x}
    }$}
\end{equation}
\begin{equation}
    M_{eq} = M + \frac{I_{2a} + I_{2b}}{r_{2a}^2} + \frac{I_{3b}\,r_{2a,cg}^2}{r_{3b,cr}^2\,r_{2a}^2}
\end{equation}

\subsubsubsection{Limits}
The following conditions must hold for the model assumptions to remain valid. Violation of any one signals a regime change: wheel lift, chain slack, or sliding contact:
\begin{align}
    {}^{z_0}R_{2a/0} &\le 0, \qquad {}^{z_0}R_{2b/0} \le 0 \qquad \text{(no wheel lift)} \\
    \left|{}^{x_0}R_{2a/0}\right| &\le -\mu_s\,{}^{z_0}R_{2a/0}, \qquad
    \left|{}^{x_0}R_{2b/0}\right| \le -\mu_s\,{}^{z_0}R_{2b/0} \qquad \text{(no slip)} \\
    T_{2a/3b} &\ge 0 \qquad \text{(no freewheel)}
\end{align}

\subsection{Analysis}

\subsubsubsection{Assumptions}
\begin{itemize}
    \item Steady state --- $\ddot{x}^{(1)} = 0$, constant speed on constant gradient
    \item No braking --- $B_{2a/1} = B_{2b/1} = 0$
    \item Forward motion --- $\sigma = 1$
    \item Still air or head/tailwind only --- drag enters as $\tfrac{1}{2}\rho\,C_dA\,v_{rel}^2$ ($\rho$ here is air density)
    \item Drivetrain efficiency not modelled --- power input equals power at the chain
    \item All limits satisfied --- no slip, no wheel lift, no freewheel
\end{itemize}

\subsubsubsection{Equilibrium}
Setting $\ddot{x}^{(1)} = 0$, $B_{2a/1} = B_{2b/1} = 0$, $\sigma = 1$ gives the steady-state force balance. Three resistance terms appear: gradient, rolling, and aerodynamic drag, scaling respectively as $v^0$, $v^0$, and $v^2$:
\begin{equation}
    \frac{r_{1,pedal}\,r_{2a,cg}}{r_{3b,cr}\,r_{2a}}\left(P_{3b/1,L} + P_{3b/1,R}\right) = Mg\left(c_{2a,rr}\cos\theta + \sin\theta\right) + \tfrac{1}{2}\rho\,C_dA\,v^2
\end{equation}
Expressing the left side as prescribed power $\mathcal{P}$ using $|\dot{\phi}_{3b}| = r_{2a,cg}v/(r_{3b,cr}r_{2a})$ eliminates the gear geometry:
\begin{equation}
    \frac{\mathcal{P}}{v} = Mg\left(c_{2a,rr}\cos\theta + \sin\theta\right) + \tfrac{1}{2}\rho\,C_dA\,v^2
\end{equation}

\subsubsubsection{Sensitivity}
Write the aerodynamic drag power as $k \equiv \tfrac{1}{2}\rho\,C_dA$ and drop rolling resistance
(an order of magnitude below gravity and drag at any sustained speed), leaving the simplest model
that still contains both effects each question is about:
\begin{equation}
    \mathcal{P} = Mg\,v\sin\theta + k\,v^3
\end{equation}

\subparagraph{Elasticities.}
The elasticity $E_p \equiv (p/v)(\partial v/\partial p)$ gives the fractional change in speed per
fractional change in parameter $p$, at constant $\mathcal{P}$. Write the balance as
$G(v;\,M,k,\theta,\mathcal{P}) = Mg\,v\sin\theta + k\,v^3 - \mathcal{P} = 0$ and
differentiate implicitly:
\begin{equation}
    \frac{\partial v}{\partial p} = -\frac{\partial G/\partial p}{\partial G/\partial v},
    \qquad
    \frac{\partial G}{\partial v} = Mg\sin\theta + 3kv^2 \equiv D
\end{equation}
\begin{equation}
    \boxed{E_M = -\frac{Mg\sin\theta}{D}, \qquad
           E_k = -\frac{kv^2}{D}, \qquad
           E_{\mathcal{P}} = \frac{Mg\sin\theta + kv^2}{D}}
\end{equation}
Since $G$ is linear in $(M, k, \mathcal{P})$ at fixed $v$, Euler's theorem gives the exact identity:
\begin{equation}
    E_M + E_k + E_{\mathcal{P}} = 0
\end{equation}

\subparagraph{Two gradient limits.}
\textbf{Flat, aero-dominated} ($Mg\sin\theta \ll kv^2$, $D \approx 3kv^2$):
\begin{equation}
    E_M \approx 0, \qquad E_k \approx -\tfrac{1}{3}, \qquad E_{\mathcal{P}} \approx \tfrac{1}{3}
    \quad\Longrightarrow\quad v \propto \mathcal{P}^{1/3},\;\; v \propto k^{-1/3}
\end{equation}
Weight is irrelevant on the flat; speed depends only on power and drag, and weakly (cube root).

\textbf{Steep, gravity-dominated} ($kv^2 \ll Mg\sin\theta$, $D \approx Mg\sin\theta$):
\begin{equation}
    E_M \approx -1, \qquad E_k \approx 0, \qquad E_{\mathcal{P}} \approx 1
    \quad\Longrightarrow\quad v \approx \frac{\mathcal{P}}{Mg\sin\theta}
\end{equation}
Aerodynamics is irrelevant while climbing; speed is set entirely by power-to-weight.

\subparagraph{Does aerodynamics matter more for a lighter rider?}
Two riders with the same $\mathcal{P}$ and $k$, differing only in mass. Since $E_M < 0$ the
lighter rider is faster. Differentiate $|E_k| = kv^2/D$ along the equilibrium curve:
\begin{equation}
    \frac{\mathrm{d}|E_k|}{\mathrm{d}v} = \frac{2kv\,Mg\sin\theta}{D^2} > 0
\end{equation}
$|E_k|$ is strictly increasing in $v$, so aerodynamics matters strictly more to the lighter
(faster) rider. The gap is smallest while climbing ($|E_k|\approx 0$ for both) and largest on
the flat ($|E_k|\to\tfrac{1}{3}$ for both).

\subparagraph{Is there a best rider weight?}
Comparing riders of different size: power and frontal area both scale with mass. Model with
exponents against a reference rider $(M_0, \mathcal{P}_0, k_0)$:
\begin{equation}
    \mathcal{P}(M) = \mathcal{P}_0\!\left(\frac{M}{M_0}\right)^{a}, \qquad
    k(M) = k_0\!\left(\frac{M}{M_0}\right)^{b}
\end{equation}
($b \approx \tfrac{2}{3}$ from geometric similarity; $a$ is physiological, typically $a < 1$.)
The size elasticity follows by chain rule:
\begin{equation}
    \boxed{\left.\frac{\mathrm{d}\ln v}{\mathrm{d}\ln M}\right|_{\text{across riders}}
    = E_M + a\,E_{\mathcal{P}} + b\,E_k}
\end{equation}
Plugging in the two limits: climbing gives $\approx a - 1 < 0$ (smaller always wins for $a < 1$);
flat gives $\approx (a-b)/3$ (bigger wins if power scales faster than frontal area, $a > b$).
Best weight flips sign with gradient; the crossover depends on $a - b$, not on mass directly.

\subparagraph{Is a steeper climb more energy-efficient for the same height gain?}
Fix height $\Delta h$; route length is $d(\theta) = \Delta h / \sin\theta$. Energy spent at
constant power $\mathcal{P}$:
\begin{equation}
    E(\theta) = \mathcal{P}\,\frac{d(\theta)}{v(\theta)}
    = \frac{\mathcal{P}\,\Delta h}{v(\theta)\sin\theta}
\end{equation}
From the power balance, $v\sin\theta = (\mathcal{P} - kv^3)/(Mg)$, which is maximised by
minimising $v$. Since $\partial v/\partial\theta < 0$ (steeper is slower),
$v\sin\theta$ rises monotonically with gradient and $E(\theta)$ falls:
\begin{equation}
    \frac{\mathrm{d}E}{\mathrm{d}\theta} < 0
    \;\Longrightarrow\;
    \text{steeper is always more energy-efficient for a fixed height gain}
\end{equation}
Mechanism: a steeper route is shorter and slower; aerodynamic loss scales as $v^2 \times
\text{distance}$, so cutting speed helps cubically while the shorter distance helps linearly.

\subsubsubsection{Stoppie limit}
Front braking transfers load forward. The rear wheel lifts when ${}^{z_0}R_{2a/0} = 0$. Setting this in the normal reaction expression fixes $K$, which is then substituted into its definition to give the critical deceleration:
\begin{equation}
    {}^{z_0}R_{2a/0} = 0 \;\Rightarrow\; K = l_D'\,Mg\cos\theta
\end{equation}
\begin{equation}
    \ddot{x}^{(1)}_{\text{stoppie}} = \frac{-\left(l_D'\,Mg\cos\theta + hM\,g\sin\theta + r_{2a}\,A_{1,x} - A_{1,m}\right)}{J/r_{2a} + hM}
\end{equation}

\subsubsubsection{Wheelie limit}
Hard acceleration transfers load rearward. The front wheel lifts when ${}^{z_0}R_{2b/0} = 0$:
\begin{equation}
    {}^{z_0}R_{2b/0} = 0 \;\Rightarrow\; K = -l_C'\,Mg\cos\theta
\end{equation}
\begin{equation}
    \ddot{x}^{(1)}_{\text{wheelie}} = \frac{l_C'\,Mg\cos\theta - hM\,g\sin\theta - r_{2a}\,A_{1,x} + A_{1,m}}{J/r_{2a} + hM}
\end{equation}

\subsubsubsection{Traction limit}
Slip is checked after the normal reactions are known. It does not feed back into the EOM --- the no-slip assumption is either valid or it is not, and violation signals that a different model (e.g.\ with combined slip) is required. The available friction at each contact is $\mu_s\left|{}^{z_0}R_{k/0}\right|$; the longitudinal force must stay within it:
\begin{align}
    \left|{}^{x_0}R_{2a/0}\right| &\le -\mu_s\,{}^{z_0}R_{2a/0} \;\Rightarrow\; \text{rear contact valid} \\
    \left|{}^{x_0}R_{2b/0}\right| &\le -\mu_s\,{}^{z_0}R_{2b/0} \;\Rightarrow\; \text{front contact valid}
\end{align}

\subsubsubsection{Todo}
\begin{itemize}
    \item FBD / dimension diagrams
\end{itemize}
\end{aiblock}


\subsection{Turning motion}
\ai{Decision: for now, turning stays Simplified-only -- no dedicated treatment in 2D/3D yet.}
\subsubsubsection{Todo}
\begin{itemize}
    \item everything
\end{itemize}

\newpage
\section{2D - Models}
\label{sec:2dmodels}
\subsection{Forces and moments}
\label{sec:2dmodels-forces}
Every external or reaction force used in 2D -- Systems belongs to one of a small number of physical types, catalogued here. This is about what each type physically represents; symbol and subscript rules are in Appendix~\ref{app:notation-forces}.

\subsubsection{Weight ($W$)}
Acts at each body's own CoG, point $a$.
\begin{equation}
    \vv{W}_k = -m_k g\,\hat{z}_i
\end{equation}
$m_k$ -- mass of body $k$, $g$ -- gravitational acceleration (\S\ref{sec:gravity}), $\hat{z}_i$ -- direction, assumed vertical.

\subsubsection{Aerodynamic ($A$)}
Drag force and pitching moment from relative airflow, quasi-steady. Treated as a known, prescribed input.
\begin{equation}
    \vv{A}_{k,f} = -A_{k,x}\,\hat{x}_1 - A_{k,z}\,\hat{z}_1, \qquad \vv{A}_{k,m} = -A_{k,m}\,\hat{y}
\end{equation}
$\hat{x}_1,\hat{z}_1$ -- frame+rider's own axes, used for every body $k$'s aerodynamic force regardless of which body it acts on. $A_{k,x}$, $A_{k,z}$, $A_{k,m}$ -- signed scalars, defined in \S\ref{sec:aerodynamics-model}.

\subsubsection{Reaction ($R$)}
Contact and joint constraint forces -- axle bearings, ground contact. Not modeled from a force law: a free, unknown planar vector, solved from the equations of motion plus the kinematic constraints in Assembly.
\begin{equation}
    \vv{R}_{k/l} = R_{k/l,x}\,\hat{x} + R_{k/l,z}\,\hat{z}
\end{equation}
The axis is resolved per instance where needed -- e.g.\ ground contact resolves along the ground-local frame $\hat{x}_0,\hat{z}_0$, with the unilateral-contact/Coulomb-friction law in \S\ref{sec:wheels-ground} written as projections onto it.

\subsubsection{Pedal force ($P$)}
Pedaling magnitude, the rider's control input -- not derived from a physical law like the others, the quantity Optimization (Timeline) eventually solves for.
\begin{equation}
    \vv{P}_{3b/1,X} = {}^{x_{3b}}P_{3b/1,X}\,\hat{x}_{3b} + {}^{z_{3b}}P_{3b/1,X}\,\hat{z}_{3b}, \qquad X\in\{D,F\}
\end{equation}
${}^{x_{3b}}P_{3b/1,X}$, ${}^{z_{3b}}P_{3b/1,X}$ -- respectively the radial and axial contributions. Both components free, for each pedal point $D$, $F$ (\S\ref{sec:frame}).

\subsubsection{Bearing friction ($Q$)}
A resisting torque at each axle, opposing relative spin between the two bodies sharing it.
\begin{equation}
    \vv{Q}_{k/l} = -|Q_{k/l}|\,\sgn(\omega_{k/l})\,\hat{y}
\end{equation}
$\omega_{k/l}$ -- relative angular velocity of body $k$ w.r.t.\ body $l$ (\S\ref{app:notation-kinematics}).


\subsection{Rigid body dynamics}
\label{sec:rigidbody}
Every body $k$ uses its CoG position $\vv{x}_k$ and (planar) rotation angle $\theta_k$ as state, with translational velocity/acceleration $\vv{v}_k,\vv{a}_k$ and angular velocity/acceleration $\omega_k,\alpha_k$ as their respective time derivatives. Since motion is planar, $I_k$ is always a scalar (moment of inertia about $\hat y$) -- no inertia tensor needed, regardless of whether a body's own mass distribution is axisymmetric:
\begin{equation}
    \sum\vv{F}_k = m_k\vv{a}_k, \qquad \sum\vv{M}_k(a_k) = I_k\alpha_k\,\hat y.
\end{equation}
Moments about a point $P\ne a_k$ pick up a transport term, from $\vv{M}_P=\vv{M}_{a_k}+\vv{r}_{a_k/P}\times\vv{F}_{total}$:
\begin{equation}
    \sum\vv{M}_k(P) = I_k\alpha_k\,\hat y + \vv{r}_{a_k/P}\times m_k\vv{a}_k.
\end{equation}
Each system's own Equations of motion just instantiates this.


\subsection{Models}
\label{sec:2dmodels-models}
Functional forms and numeric values behind the prescribed forces above.

\subsubsection{Mass, moment of inertia, and center of gravity}
\label{sec:mass-inertia}
Each body's mass $m_k$, planar moment of inertia $I_k$ (about $\hat y$, through its own CoG), and CoG location (point $a$, \S\ref{app:notation-points}) are physical inputs, not modeled from a law -- measured (CAD, real life) or otherwise assumed, feeding Weight (\S\ref{sec:2dmodels-forces}) and Rigid body dynamics (\S\ref{sec:rigidbody}). For now $I_k$ and CoG location are kept constant for every body; $m_k$ is constant too, except for the frame+rider system, where it varies (rider mass) -- detailed in \S\ref{sec:frame}.

\subsubsection{Gravity}
\label{sec:gravity}
Uniform field, constant magnitude $g$ (standard value $9.81\,\text{m/s}^2$ unless overridden). Real-world $g$ varies with latitude and altitude -- roughly $9.78\,\text{m/s}^2$ at the equator to $9.83\,\text{m/s}^2$ at the poles (sea level), a spread under $0.5\%$ across any cyclable terrain -- small enough that constant $g$ is the default. Where altitude dependence is wanted instead, the inverse-square model
\begin{equation}
    g(h) = g_0\left(\frac{r_E}{r_E+h}\right)^2
\end{equation}
is available, with $g_0$ the standard surface value, $r_E$ Earth's radius, and $h$ altitude above sea level; latitude still isn't captured by it.

\subsubsection{Aerodynamics}
\label{sec:aerodynamics-model}
Drag, side force, and pitching moment, quasi-steady, each a dynamic-pressure-scaled coefficient along the frame's own axes $\hat{x}_1,\hat{z}_1,\hat y$ (\S\ref{sec:2dmodels-forces}):
\begin{align}
    A_{k,x} &= q\,C_{k,x}\,A_k, \\
    A_{k,z} &= q\,C_{k,z}\,A_k, \\
    A_{k,m} &= q\,C_{k,m}\,A_k\,l_k.
\end{align}

\textbf{Dynamic pressure.}
\begin{equation}
    q = \tfrac{1}{2}\rho\,v_{rel}^2,
\end{equation}
with $\rho$ air density (\S\ref{sec:air-density}) and $v_{rel}=|\vv{v}_{rel}|$ the airspeed, defined next.

\textbf{Relative velocity.}
\begin{equation}
    \vv{v}_{rel} = \vv{v}_{wind} - \vv{v}_1,
\end{equation}
the air's velocity relative to the frame+rider system.

\textbf{Coefficients.} Each is empirical, generally only available as a wind-tunnel or CFD lookup table over its own arguments:
\begin{equation}
    C_{k,x},\,C_{k,z},\,C_{k,m} = C(\xi_k,\,\omega_k,\,\beta_{rel},\,Re)
\end{equation}
where $\xi_k$ is a geometric/configuration change of the body itself (e.g.\ rider posture for $k=1$; rigid single-piece bodies like the wheels have none), $\omega_k$ the body's own rotation rate (e.g.\ wheel spin), $\beta_{rel}$ the angle of attack, the relative wind's angle to the body (\S\ref{app:notation-kinematics}):
\begin{equation}
    \beta_{rel} = \operatorname{atan2}\!\big({}^{z_1}v_{rel},\,{}^{x_1}v_{rel}\big),
\end{equation}
and $Re$ Reynolds number:
\begin{equation}
    Re = \frac{\rho\,v_{rel}\,l_k}{\mu},
\end{equation}
with $\mu$ air's dynamic viscosity and $l_k$ the reference length below.

For simplicity, each coefficient can instead be linearized -- a first-order expansion in these arguments around a reference point:
\begin{equation}
    C = C_0 + \frac{\partial C}{\partial \xi_k}\,\xi_k + \frac{\partial C}{\partial \omega_k}\,\omega_k + \frac{\partial C}{\partial \beta_{rel}}\,\beta_{rel} + \frac{\partial C}{\partial Re}\,Re,
\end{equation}
written generically for any of $C_{k,x},C_{k,z},C_{k,m}$, each with its own constant $C_0$ and four slopes. The $\beta_{rel}=0$ intercept $C_{k,x,0}$ recovers exactly the lumped $C_dA$ Timeline's Data acquisition calibrates from ordinary ride data.

\textbf{Reference area and length.} $A_k$ is body $k$'s frontal reference area (projected silhouette normal to $\hat x_1$) and $l_k$ a reference length for the moment arm (e.g.\ wheelbase) -- fixed geometric inputs, same status as \S\ref{sec:mass-inertia}'s mass/inertia/CoG: measured (CAD, real life) or otherwise assumed.

\subsubsection{Atmosphere}
Atmospheric data -- density, temperature, pressure, and humidity -- can be assumed constant, or obtained from a model:
\begin{itemize}
    \item \textbf{ISA} (International Standard Atmosphere) -- standard altitude-based model; see Appendix~\ref{app:isa}.
    \item \textbf{Weather station data} -- a single station, or multiple stations interpolated along the route.
\end{itemize}

\subsubsubsection{Air density}
\label{sec:air-density}
\begin{equation}
    \rho = \rho(h, RH)
\end{equation}
Altitude dependence comes from the ISA model (\S\ref{app:isa}), dry air only, via the ideal-gas relation given there. Real (wet) air density additionally needs relative humidity $RH$ from weather-station data, folded in through the moist-air correction (\S\ref{app:humidity-density}).

\subsubsubsection{Viscosity}
\label{sec:air-viscosity}
\begin{equation}
    \mu = \mu(T)
\end{equation}
$\mu$ (\S\ref{sec:aerodynamics-model}'s $Re$) depends primarily on temperature -- rising with it, unlike a liquid's viscosity, which falls; see Sutherland's law (\S\ref{app:sutherland}) for the model.

\subsubsection{Pedal force}
\label{sec:pedal-force-model}
Not modeled from a force law either: $\vv{P}_{3b/1,X}$ (\S\ref{sec:2dmodels-forces}) is the rider's control input, free within physical pedaling limits, solved for by Optimization rather than prescribed here (Timeline).

\textbf{Periodic analytical.} Where a closed-form model is wanted instead of a free decision variable, each pedal's magnitude can be modeled as a periodic function of crank angle $\theta_{3b}$, the two pedals offset by half a revolution:
\begin{equation}
    P_D(\theta_{3b}) = P_0\cos^2(\theta_{3b}-\phi), \qquad P_F(\theta_{3b}) = P_0\cos^2(\theta_{3b}-\phi+\pi),
\end{equation}
squared rather than left as a bare cosine so the magnitude never goes negative, with $P_0$ peak magnitude and $\phi$ the crank angle of peak force. The same squared-for-non-negativity structure extends to more harmonics:
\begin{equation}
    P_X(\theta_{3b}) = \left(a_0 + \sum_{n=1}^N a_n\cos(n\theta_{3b}-\phi_n)\right)^2, \qquad X\in\{D,F\},
\end{equation}
$F$ offset from $D$ by $\pi$ in $\theta_{3b}$ as above, with coefficients $a_0,\dots,a_N,\phi_1,\dots,\phi_N$ free to fit.

\textbf{Shape-constrained templates.} Real pedaling isn't as symmetric as $\cos^2$ suggests -- tangential force collapses near top and bottom dead center ($\theta_{3b}=0,\pi$) and peaks around $90^\circ$--$110^\circ$ past top dead center. Two analytical templates capture this with a couple of shape parameters, $X=F$ shifted from $X=D$ by $\pi$ in $\theta_{3b}$ as above:
\begin{align}
    P_X(\theta_{3b}) &= P_0\,\big[\max(0,\sin(\theta_{3b}-\theta_0))\big]^n, \\
    P_X(\theta_{3b}) &= P_0\,\exp\!\big(\kappa(\cos(\theta_{3b}-\theta_0)-1)\big).
\end{align}
The first is a powered, dead-zone-clipped sine: raising $n$ sharpens the peak and pushes the shoulders to exactly zero over a wider span around dead center, instead of $\cos^2$'s soft, never-quite-zero approach. The second is a von Mises-type circular bump -- the natural ``circular Gaussian'' for angle-indexed data -- where concentration $\kappa$ controls peak sharpness, and it's smooth and non-negative everywhere without a hard clip. Both reduce to something close to $\cos^2$ as the shape parameter relaxes ($n=2$ or $\kappa\to0$). $P_0$ is peak magnitude and $\theta_0$ the crank angle of peak force in both.

\textbf{Biomechanical (legs isolation).} The deepest option: a multi-segment leg model (hip--knee--ankle) with muscle actuators, mapped through the pedal/crank kinematic chain to the resultant pedal force -- predicts the force profile from physiology rather than fitting a shape to it. Deferred to the rider legs-isolation model (\S\ref{sec:frame}'s pending ``rider legs isolated'' assumption).

\textbf{Data-driven.} Skip the functional form entirely and use measured pedal-force or power-meter data directly, per rider, once available (Timeline's Data acquisition) -- either fit it to whichever of the above forms is chosen, or use a normalized torque-curve shape (own data, or from cycling-biomechanics literature) as a fixed template, scaled by rider-specific peak force and cadence.

\subsubsection{Bearing friction}
\label{sec:bearing-friction-model}
Only the resisting moment $\vv{Q}_{k/l}$ (\S\ref{sec:2dmodels-forces}) is modeled for now, no bearing friction force. Standard rolling-element bearing friction (Coulomb, load-proportional):
\begin{equation}
    |Q_{k/l}| = \tfrac{1}{2}\,\mu_b\,d_{k/l}\,|\vv{R}_{k/l}|,
\end{equation}
with $\mu_b$ an empirical bearing-type friction coefficient (e.g.\ $\approx0.0015$ for deep-groove ball bearings), $d_{k/l}$ the bearing bore/shaft diameter at that joint, and $|\vv{R}_{k/l}|$ the reaction force magnitude the bearing carries (\S\ref{sec:2dmodels-forces}).


\newpage
\section{2D - Systems}
\label{sec:2dsystems}
\subsubsection{Systems}
\begin{enumerate}
    \item Frame + Rider
    \item Wheel
        \begin{enumerate}
            \item Rear
            \item Front
        \end{enumerate}
    \item Drivetrain
        \begin{enumerate}
            \item Casette
            \item Crankset
        \end{enumerate}
\end{enumerate}

\subsubsection{Current global assumptions}
\begin{itemize}
    \item Planar (sagittal-plane) motion only: every rigid body has 3 DOF -- translation along $\hat x_i,\hat z_i$ and rotation about $\hat y$. No roll, yaw, or out-of-plane motion.
    \item All bodies (frame+rider, wheels, drivetrain components) are rigid: no structural flex, no tire deformation. \emph{Pending: frame equivalent stiffness, rim/wheel stiffness, tire stiffness -- see Appendix~\ref{app:refinements}.}
    \item Frame $i$ is the fixed, inertial (non-accelerating, non-rotating) world frame; $\hat z_i$ points opposite gravity, gravitational acceleration $g$ is constant. Frame $0$ is the ground-local frame, tied to the road at the current contact point (matches the ground-body subscript in $\vv{R}_{2a/0}$).
    \item The road is a known, smooth ($C^1$, ideally $C^2$) graph $z(x)$ in the $x_i$-$z_i$ plane -- matches how real elevation data (GPS routes, surveys, DEMs) is normally given. Single-valued: no overhangs or vertical loops, which is fine for road terrain but will need generalizing to an arc-length curve for BMX/MTB features later. See Assembly for the frame-$0$ construction.
    \item Aerodynamic loads ($\vv{A}$) are treated as known/prescribed quasi-steady inputs, not derived from a flow model.
    \item Mass, CoG location, and inertia of each rigid body are known constants.
\end{itemize}

\subsubsection{Notation}
Dimension, force/moment, and point labeling conventions used throughout this document are collected in Appendix~\ref{app:notation}.


\newpage
\subsection{Frame + Rider}
\label{sec:frame}
Single instance -- no rear/front split like Wheel or Drivetrain, so the subsystem-instance heading level is skipped here and this goes straight to the content categories.
\subsubsubsection{Assumptions}
\begin{itemize}
    \item Rider is rigidly attached to the frame; lumped as one rigid body (no relative rider motion, e.g. standing or leaning, modeled yet). \emph{Pending: rider legs isolated, rider fully isolated.}
    \item CoG location ($l_{1,a},h_{1,a}$) and points $B$, $C$ are fixed in the frame's body frame ($\hat x_1,\hat z_1$): no suspension travel, rigid fork and rear triangle. \emph{Pending: suspension (front fork and rear).}
    \item Pedaling force is applied at fixed radius $r_{1,pedal}$ about $C$ at the crank's current orientation $\hat x_{2b}$ -- an instantaneous snapshot, not yet coupled to crank rotation dynamics. \emph{Pending: complex drivetrain.}
\end{itemize}

\subsubsubsection{Force}
The equation of motion for this system is given by
\begin{equation}
    \sum \vv{F}_1 = \vv{W}_1 + \vv{A}_{1,f} + \vv{R}_{1/2a} + \vv{R}_{1/2b} + \vv{R}_{1/3b} + \vv{P}_{1/3b, D} + \vv{P}_{1/3b, F},
\end{equation}
where respectively the forces acting on the system are weight, aerodynamic force, wheel reactions and drivetrain reactions (including bearing and pedaling force).

\subsubsubsection{Geometry}
The geometry will be derived with respect to point $A$ at center of rear axle, where aerodynamics and rear wheel reaction forces act.

Weight acts at point $a$ center of gravity (CoG)
\begin{equation}
    \vv{Aa} = l_{1,a}\hat{x}_1 + h_{1,a}\hat{z}_1.
\end{equation}

Front wheel reaction acts at $B$
\begin{equation}
    \vv{AB} = l_{1,B}\hat{x}_1.
\end{equation}

The drivetrain bearing reaction acts at $C$
\begin{equation}
    \vv{AC} = l_{1,C}\hat{x}_1 + h_{1,C}\hat{z}_1,
\end{equation}
while pedaling forces act radially around $C$ in $F$ and $B$, giving
\begin{align}
    \vv{AD} &= \vv{AC} + \vv{CD} = l_{1,C}\hat{x}_1 + h_{1,C}\hat{z}_1 + r_{1,pedal} \hat{x}_{2b}, \\
    \vv{AF} &= \vv{AC} + \vv{CF} = l_{1,C}\hat{x}_1 + h_{1,C}\hat{z}_1 - r_{1,pedal} \hat{x}_{2b}.
\end{align}

\subsubsubsection{Moment}
Rotation equation of motion are given as
\begin{align}
    \sum \vv{M}_1(A) &= \vv{Aa} \times \vv{W}_1 + \vv{A}_{1,m} + \vv{AB} \times \vv{R}_{1/2b} + \vv{AC} \times \vv{R}_{1/3b} + \vv{AD} \times \vv{P}_{1/3b, D} + \vv{AF} \times \vv{P}_{1/3b, F}, \\
    &+ \vv{Q}_{1/2a} + \vv{Q}_{1/2b} + \vv{Q}_{1/3b},
\end{align}
where $\vv{A}_{1,m}$ is aerodynamic pitching moment and $Q$ are friction moments at bearings.


\subsubsubsection{Equations of motion}
\begin{align}
    \sum\vv{F}_1 &= m_1\vv{a}_1, \\
    \sum\vv{M}_1(A) &= I_1\alpha_1\,\hat y + \vv{Aa}\times m_1\vv{a}_1.
\end{align}

 \newpage
\subsection{Wheel}
\label{sec:wheel}

Common assumptions (both instances):
\begin{itemize}
    \item Circular, axisymmetric rigid disk of constant radius $r$, assumed equal both instances ($r_{2a}=r_{2b}=r$): no tire deformation, no out-of-round effects. \emph{Pending: tire stiffness, rim/wheel stiffness.}
    \item CoG offset from the hub, $r_a$ ($r_{2a,a}$, $r_{2b,a}$), is kept general to allow imbalance between the two wheels.
    \item Ground contact idealized as a single point, valid while the local road radius of curvature is much larger than $r$ -- see Assembly. \emph{Pending: tire stiffness (a deformable tire has a finite contact patch, not a point).}
\end{itemize}

\subsubsection{Rear (2a)}
Carries the drivetrain (cassette) reaction at the hub.

\subsubsubsection{Force}
The equation of motion for this system is given by
\begin{equation}
    \sum \vv{F}_{2a} = \vv{W}_{2a} + \vv{A}_{2a,f} + \vv{R}_{2a/0} + \vv{R}_{2a/1} + \vv{R}_{2a/3a, f},
\end{equation}
where respectively the forces acting on the system are weight, aerodynamic force, ground reaction, frame reaction and drivetrain reaction.

\subsubsubsection{Geometry}
The geometry will be derived with respect to point $A$ at center of axle, where aerodynamics, frame reaction and drivetrain reaction force act at.

Weight acts at point $a$ center of gravity (CoG)
\begin{equation}
    \vv{Aa} = r_{2a,a}\hat{x}_{2a}
\end{equation}

Ground reaction acts at $b$,
\begin{equation}
    \vv{Ab} = l_{2a,rr}\,\sigma\,\hat x_0(x^*) - r_{2a}\,\hat z_0(x^*).
\end{equation}

\subsubsubsection{Moment}
Rotation equation of motion are given as
\begin{align}
    \sum \vv{M}_{2a}(A) &= \vv{Aa} \times \vv{W}_{2a} + \vv{A}_{2a,m} + \vv{Ab} \times \vv{R}_{2a/0} + \vv{R}_{2a/3a, m} + \vv{Q}_{2a/1},
\end{align}
where $\vv{R}_{2a/3a, m}$ is drivetrain fixed reaction force.

\subsubsubsection{Equations of motion}
\begin{align}
    \sum\vv{F}_{2a} &= m_{2a}\vv{a}_{2a}, \\
    \sum\vv{M}_{2a}(A) &= I_{2a}\alpha_{2a}\,\hat y + \vv{Aa}\times m_{2a}\vv{a}_{2a}.
\end{align}

\subsubsection{Front (2b)}
No drivetrain reaction at the hub (no cassette on the front wheel); otherwise identical to Rear (2a).

\subsubsubsection{Force}
\begin{equation}
    \sum \vv{F}_{2b} = \vv{W}_{2b} + \vv{A}_{2b,f} + \vv{R}_{2b/0} + \vv{R}_{2b/1},
\end{equation}
i.e. the Rear (2a) force equation with the drivetrain-reaction term $\vv{R}_{2a/3a,f}$ dropped.

\subsubsubsection{Geometry}
Identical in form to Rear (2a) with $2a\to2b$, \textbf{$A\to B$}.
\begin{equation}
    \vv{Ba} = r_{2b,a}\hat{x}_{2b},
\end{equation}
and ground reaction acting at $b$,
\begin{equation}
    \vv{Bb} = l_{2b,rr}\,\sigma\,\hat x_0(x^*) - r_{2b}\,\hat z_0(x^*).
\end{equation}

\subsubsubsection{Moment}
Taken about $B$ (this wheel's own axle point, not $A$):
\begin{equation}
    \sum \vv{M}_{2b}(B) = \vv{Ba} \times \vv{W}_{2b} + \vv{A}_{2b,m} + \vv{Bb} \times \vv{R}_{2b/0} + \vv{Q}_{2b/1},
\end{equation}
i.e. the Rear (2a) moment equation with the drivetrain fixed-reaction term $\vv{R}_{2a/3a,m}$ dropped.

\subsubsubsection{Equations of motion}
\begin{align}
    \sum\vv{F}_{2b} &= m_{2b}\vv{a}_{2b}, \\
    \sum\vv{M}_{2b}(B) &= I_{2b}\alpha_{2b}\,\hat y + \vv{Ba}\times m_{2b}\vv{a}_{2b}.
\end{align}
 \newpage
\subsection{Drivetrain}

\newpage
\section{2D - Assembly}
\subsection{Frame -- Wheels}
\label{sec:frame-wheels}
Both axles are revolute (pin) joints: frame points $A$, $B$ and each wheel's own axle center, labeled the same (see \S\ref{app:notation-points}), are the same material points at all times,
\begin{equation}
    \vv{A_1A_{2a}} = \vv{0}, \qquad \vv{B_1B_{2b}} = \vv{0} \qquad \forall t,
\end{equation}
hence the material points' velocities and accelerations match too, $\vv{v}_{1,A}=\vv{v}_{2a,A}$, $\vv{a}_{1,A}=\vv{a}_{2a,A}$ (and the $B$ analogs).

Each wheel spins freely about the hub axis $\hat y$, by the relative spin angle $\phi_{2a}$ (and $\phi_{2b}$) -- the quantity Drivetrain will ultimately drive, since gearing relates relative, not absolute, rotation rates:
\begin{align}
    \theta_{2a} &= \theta_1 + \phi_{2a}, \\
    \theta_{2b} &= \theta_1 + \phi_{2b}.
\end{align}

\subsection{Wheels -- Ground}
\label{sec:wheels-ground}

\subsubsection{Contact point}
The road vector is given by
\begin{equation}
    \vv{Ob_{2a}}(x) = x\,\hat x_i + z(x)\hat z_i,
\end{equation}
where $O$ is the origin of frame $i$ and $b_{2a}$ is wheel 2a's candidate contact point on the road at horizontal position $x$ (matching the ground-contact point label $b$ already used in \S\ref{sec:wheel}) -- real elevation data (GPS routes, road surveys, digital elevation models) is almost always given this way already, as a graph $z(x)$ rather than a general parametrized curve.

With local slope angle $\theta(x) = \arctan\!\big(z'(x)\big)$, at each point $x$ define the ground-local frame $0$, tilted from frame $i$ by $\theta(x)$,
\begin{align}
    \hat x_0(x) &= \cos\big(\theta(x)\big)\,\hat x_i + \sin\big(\theta(x)\big)\,\hat z_i, \\
    \hat z_0(x) &= -\sin\big(\theta(x)\big)\,\hat x_i + \cos\big(\theta(x)\big)\,\hat z_i.
\end{align}

For wheel 2a (center $A$, radius $r_{2a}$), the geometric contact reference $x^*$ is the point on the road tangent to the wheel: chosen so the vector from the road to $A$ is normal to the road there,
\begin{equation}
    \vv{b_{2a}A}(x^*) \cdot \hat x_0(x^*) = 0.
\end{equation}

The actual contact resultant is not purely radial from $x^*$. Under load, tire deformation shifts the resultant slightly forward of $x^*$, in the direction of travel. Rather than add a separate resistance moment, that shift is folded directly into where the contact vector points:
\begin{equation}
    \vv{Ab_{2a}} = l_{2a,rr}\,\sigma\,\hat x_0 - r_{2a}\,\hat z_0, \qquad l_{2a,rr} = c_{2a,rr}\,r_{2a}, \qquad \sigma = \sgn\!\big(\vv{v}_{2a,A}\cdot\hat x_0\big),
\end{equation}
with $c_{2a,rr}$ the rolling resistance coefficient and $\sigma$ keeping the shift pointing in the current direction of travel (frame $0$ here is evaluated at $x^*$, per the previous equation).

\subsubsection{Contact reaction}
Define the gap
\begin{equation}
    g_N = \vv{b_{2a}A}\cdot\hat z_0(x^*) - r_{2a},
\end{equation}
which splits the reaction into two cases:

\textbf{Not touching ($g_N>0$):}
\begin{equation}
    \vv{R}_{2a/0} = \vv{0},
\end{equation}
no force at all.

\textbf{Touching ($g_N=0$):} $\vv{R}_{2a/0}$ is present, its components along the ground-local frame (\S\ref{sec:2dmodels-forces}) bounded by
\begin{equation}
    {}^{z_0}R_{2a/0} \ge 0, \qquad \big|{}^{x_0}R_{2a/0}\big| \le \mu_s\,{}^{z_0}R_{2a/0},
\end{equation}
the normal component non-negative (the ground can push, not pull) and the tangential component bounded by Coulomb static friction.

Together, these two cases are the unilateral-contact complementarity condition
\begin{equation}
    {}^{z_0}R_{2a/0}\ge0, \qquad g_N\ge0, \qquad {}^{z_0}R_{2a/0}\,g_N=0.
\end{equation}

\subsubsection{Rolling constraint}
While in contact, the material point $b_{2a}$ on the wheel has velocity
\begin{equation}
    \vv{v}_{2a,b_{2a}} = \vv{v}_{2a,A} + \vv{\omega}_{2a}\times\vv{Ab_{2a}},
\end{equation}
with tangential (slip) component $v_{slip} = \vv{v}_{2a,b_{2a}}\cdot\hat x_0$ (really ${}^{x_0}v_{2a,b_{2a}}$ under \S\ref{app:notation-kinematics}'s template, kept as $v_{slip}$ for readability) -- the normal component is zero automatically, enforced by the non-penetration condition above.

\textbf{Stick (no slip):} $v_{slip}=0$ is imposed as a kinematic constraint. Dropping the $l_{2a,rr}$ term from $\vv{Ab_{2a}}$ here (it is much smaller than $r_{2a}$, so its contribution to $\vv{\omega}_{2a}\times\vv{Ab_{2a}}$ is a second-order correction to the constraint, not to the moment balance where it was kept):
\begin{equation}
    {}^{x_0}v_{2a,A} = -r_{2a}\,\omega_{2a},
\end{equation}
and ${}^{x_0}R_{2a/0}$ becomes an unknown constraint force, solved from the equations of motion, valid as long as it stays within the friction bound $|{}^{x_0}R_{2a/0}|\le\mu_s\,{}^{z_0}R_{2a/0}$ above.

\textbf{Slip:} once holding $v_{slip}=0$ would require $|{}^{x_0}R_{2a/0}|>\mu_s\,{}^{z_0}R_{2a/0}$, the constraint is dropped -- $v_{slip}$ becomes a free kinematic variable governed by the dynamics -- and friction switches to the kinetic law
\begin{equation}
    {}^{x_0}R_{2a/0} = -\mu_k\,{}^{z_0}R_{2a/0}\,\sgn(v_{slip}), \qquad \mu_k \le \mu_s.
\end{equation}

Together, the contact/airborne switch and the stick/slip switch make this a hybrid (mode-switching) system; a numerical solver would typically resolve both via complementarity (LCP) or event-driven mode detection rather than a single closed-form equation set.

Same relations hold for the front wheel with $2a\to2b$, $A\to B$ (its own axle center, matching the frame's own label for the front axle -- see \S\ref{app:notation-points}), $b_{2a}\to b_{2b}$, and $r_{2a}\to r_{2b}$, $l_{2a,rr}\to l_{2b,rr}$ -- none of these need be equal front to rear, e.g.\ for a bike with different wheel sizes or tire pressures.


\newpage
\section{3D}
\label{sec:3d}
\begin{aiblock}
Final stage, much later. Full spatial (6 DOF per body) rigid-body dynamics, extending the 2D derivation. Initially scoped to motion only -- the other three Timeline tracks (optimization/strategy, R\&D, data acquisition) are only revisited here once they're finished on top of 2D. Not started.
\end{aiblock}

\newpage
\appendix
\section{Notation conventions}
\label{app:notation}
Dimension, force/moment, and point labeling conventions used throughout this document.
\todo{scan the whole document for consistency with the conventions below -- older sections may predate a rule change here and not yet match it.}
\todo{scan the whole document, including this appendix, for self-referential or history-narrating sentences (prose that explains how/why something changed rather than just stating the current rule) and remove them.}

\subsection{Forces and moments}
\label{app:notation-forces}
This is about the grammar of a force/moment symbol -- how one is built -- not what any particular letter physically means (that's \S\ref{sec:2dmodels-forces}). Every symbol follows one template:
\begin{equation}
    {}^{\text{axis}}(\text{type})_{\text{on}/\text{from},\ \text{variant},\ \text{point}}
\end{equation}

\begin{itemize}
    \item \textbf{type} -- the force/moment type letter (e.g.\ $R$, $W$, $A$, $Q$, $P$; meanings in \S\ref{sec:2dmodels-forces}). Always required.
    \item \textbf{on} -- the body the force acts \emph{on}. Always required.
    \item \textbf{from} -- the body the force acts \emph{from}. Required for a reaction pair (two modeled bodies exchanging force); omitted for a one-way/external load with no modeled ``other body'' -- gravity ($\vv{W}_1$), aerodynamics ($\vv{A}_{1,f}$).
    \item \textbf{variant} -- distinguishes multiple different forces/moments between the \emph{same} on/from pair, when needed. E.g.\ $\vv{A}_{1,f}$ (force) vs.\ $\vv{A}_{1,m}$ (moment): same interaction, two different physical quantities from it.
    \item \textbf{point} -- the point of application, when needed to distinguish multiple points for the same on/from/variant. E.g.\ $\vv{P}_{1/3b,D}$ vs.\ $\vv{P}_{1/3b,F}$: same pedal-force interaction, two different pedal positions.
    \item \textbf{axis} -- a leading superscript, not a subscript, used only when taking a scalar component of the vector along a stated axis, e.g.\ ${}^{x_0}R_{2a/0}$ for the tangential component of the rear ground reaction in the ground-local frame. It gets its own superscript slot precisely so it never collides with variant or point.
\end{itemize}

Newton's third law reads directly off an on/from pair: $\vv{R}_{k/l} = -\vv{R}_{l/k}$.

One documented exception, kept for readability rather than forced through the template: $M$ (the moment-balance equation, $\sum\vv{M}_i(X)$) isn't a physical force at all -- it takes the reference point as a function argument rather than fitting the on/from pattern.

\subsection{Displacement, velocity, acceleration}
\label{app:notation-kinematics}
Same grammar as forces and moments (\S\ref{app:notation-forces}), since a velocity needs the same three pieces of information as a reaction force: which body's point is moving, what it's measured relative to, and which point:
\begin{equation}
    {}^{\text{axis}}(\text{type})_{\text{on}/\text{from},\ \text{point}}
\end{equation}
\begin{itemize}
    \item \textbf{type} -- $x$ (displacement/position), $v$ (velocity), $a$ (acceleration); $\theta,\omega,\alpha$ for the angular counterparts (no on/from/point -- a body's own rotation isn't attached to a point).
    \item \textbf{on} -- the body whose point's motion is described. Always required.
    \item \textbf{from} -- the body the motion is measured relative to. Omitted when relative to the fixed frame $i$ (the default, exactly like a force's omitted \textbf{from} for a one-way/external load); included for relative motion between two modeled bodies.
    \item \textbf{point} -- omitted when the point is the body's own CoG, the default state point, e.g.\ $\vv{v}_i$ (\S\ref{sec:rigidbody}); included otherwise, e.g.\ $\vv{v}_{2a,A}$ for the velocity of axle point $A$ on wheel 2a.
    \item \textbf{axis} -- same leading-superscript scalar-component convention as forces, e.g.\ ${}^{x_0}v_{2a,A}$.
\end{itemize}
One exception, kept for readability: $v_{slip}$ (\S\ref{sec:wheels-ground}) is really ${}^{x_0}v_{2a,b_{2a}}$ under the full template.

\subsection{Dimensions}
Every dimension symbol follows:
\begin{equation}
    (\text{quantity})_{\text{body},\ \text{descriptor}}
\end{equation}
\begin{itemize}
    \item \textbf{quantity} -- the letter, reserved by physical meaning, not by system (table below). Always required.
    \item \textbf{body} -- which instance the value belongs to. Always required, even when two instances happen to share a value -- e.g.\ both wheels are assumed to share one radius, but it's still written $r_{2a}$, $r_{2b}$ (with $r_{2a}=r_{2b}=r$ stated separately as the assumption), not bare $r$ in an equation. Unlike points below, a dimension's body is never omitted, since there's no home-context prose paragraph to carry that information the way a point's local section does.
    \item \textbf{descriptor} -- a point letter or a name, added only when needed to tell apart multiple same-letter quantities on the same body, e.g.\ $r_{1,pedal}$ vs.\ $r_{2a,a}$, or $l_{2a,rr}$. Omitted when a body has only one quantity of that letter, e.g.\ $r_{2a}$ (wheel radius, nothing else on that body is called $r$).
\end{itemize}

The quantity letter is reserved by physical meaning, not by system, so the same letter always means the same kind of thing everywhere in the document:
\begin{center}
\begin{tabular}{c l l}
    Symbol & Meaning & Example \\
    \hline
    $l$ & body-frame in-plane (x-axis) offset & $l_{1,C}$, $l_{2a,rr}$ \\
    $h$ & body-frame vertical (z-axis) offset & $h_{1,C}$ \\
    $r$ & fixed radius, intrinsic to one body (rotation-axis distance) & $r_{2a}$, $r_{1,pedal}$, $r_{2a,a}$ \\
    $d$ & generic point-to-point distance, not axis-aligned (reserved, unused yet) & -- \\
    $\rho$ & radius of curvature of a path (extrinsic, varies with position) & $\rho(s)$
\end{tabular}
\end{center}

\subsection{Points}
\label{app:notation-points}
Simplest of the four grammars -- a point is just a letter, optionally tagged with the body it belongs to:
\begin{equation}
    (\text{letter})_{\text{body}}
\end{equation}
\begin{itemize}
    \item \textbf{letter} -- the point label itself. Its meaning is local to the system/file that defines it (home context) -- the same letter can mean a different point in a different system, same as a dimension's quantity-letter does \emph{not} work (dimensions are reserved globally; points are not). Always required.
    \item \textbf{body} -- unlike a dimension, \emph{omitted} within home context (the enclosing section already disambiguates which point is meant); attached as a subscript only when the point crosses into Assembly or a cross-system equation, e.g.\ $A_1$ for frame point $A$, $A_{2a}$ for wheel 2a's own axle point. No superscript (no $A^{(1)}$) -- the same mechanism already used for a position vector of a point, e.g.\ $\vv{x}_{1,A}$ (\S\ref{app:notation-kinematics}).
\end{itemize}

Case carries meaning: uppercase ($A$, $B$, \dots) marks a repeatable point -- one that can be reused across bodies when it genuinely coincides with another (the axle joints, \S\ref{sec:frame-wheels}). Lowercase marks a point intrinsic to one body, never repeated -- e.g.\ $a$ is reserved for the point where weight acts (CoG).


\section{Atmosphere reference}
\label{app:atmosphere}
\subsection{ISA (International Standard Atmosphere)}
\label{app:isa}
Troposphere ($h\le11\,\text{km}$, well above any cyclable altitude): temperature falls linearly with altitude, and pressure/density follow from the hydrostatic and ideal-gas relations:
\begin{align}
    T(h) &= T_0 - L h, \\
    p(h) &= p_0\left(1 - \frac{Lh}{T_0}\right)^{\textstyle\frac{g_0 M}{RL}}, \\
    \rho(h) &= \frac{p(h)}{R_d\,T(h)},
\end{align}
with sea-level reference $T_0=288.15\,\text{K}$, $p_0=101325\,\text{Pa}$, lapse rate $L=0.0065\,\text{K/m}$, standard gravity $g_0=9.80665\,\text{m/s}^2$, molar mass of air $M=0.0289644\,\text{kg/mol}$, universal gas constant $R=8.31446\,\text{J/(mol\,K)}$, and specific gas constant for dry air $R_d=R/M\approx287.05\,\text{J/(kg\,K)}$.

\begin{center}
\begin{tabular}{c c c c}
    $h$ (m) & $T$ ($^\circ$C) & $p$ (kPa) & $\rho$ (kg/m$^3$) \\
    \hline
    0 & 15.0 & 101.3 & 1.225 \\
    1000 & 8.5 & 89.9 & 1.112 \\
    2000 & 2.0 & 79.5 & 1.007 \\
    3000 & -4.5 & 70.1 & 0.909 \\
    4000 & -11.0 & 61.6 & 0.819 \\
    5000 & -17.5 & 54.0 & 0.736 \\
\end{tabular}
\end{center}

\subsection{Humidity-adjusted density}
\label{app:humidity-density}
Moist air is less dense than dry air at the same temperature and pressure (water vapor molecules are lighter than N$_2$/O$_2$), so $\rho$ from \S\ref{app:isa} needs a correction once relative humidity $RH$ is known:
\begin{equation}
    \rho = \frac{p}{R_d T}\left(1 - \Big(1-\frac{R_d}{R_v}\Big)\frac{e}{p}\right), \qquad e = RH\,e_s(T), \qquad e_s(T) = 0.61078\exp\!\left(\frac{17.27\,T}{T+237.3}\right),
\end{equation}
with $R_d\approx287.05\,\text{J/(kg\,K)}$ (\S\ref{app:isa}) and $R_v\approx461.5\,\text{J/(kg\,K)}$ the specific gas constants for dry air and water vapor respectively, $e$ the water vapor partial pressure, $RH\in[0,1]$ relative humidity, and $e_s(T)$ the saturation vapor pressure (Tetens' formula, $T$ in $^\circ$C, $e_s$ in kPa).

\subsection{Viscosity (Sutherland's law)}
\label{app:sutherland}
Dynamic viscosity of a gas rises with temperature (the opposite of a liquid), via
\begin{equation}
    \mu(T) = \mu_0\left(\frac{T}{T_0}\right)^{3/2}\frac{T_0+S}{T+S},
\end{equation}
with reference viscosity $\mu_0=1.716\times10^{-5}\,\text{Pa\,s}$ at $T_0=273.15\,\text{K}$, and Sutherland's constant for air $S=110.4\,\text{K}$.


\section{Candidate refinements}
\label{app:refinements}
\begin{aiblock}
Model refinements worth doing at some point, not yet prioritized against the 2D build order in the Timeline. Pull an item into that ordered list when it becomes the next thing to build.
\begin{itemize}
    \item \textbf{Frame equivalent stiffness} -- frame as a rigid body plus an equivalent spring-damper rather than perfectly rigid; matters most for standing/sprint efforts.
    \item \textbf{Rim/wheel stiffness} -- separate from tire stiffness.
    \item \textbf{Chain and drivetrain compliance} -- a real chain isn't infinitely stiff under load.
    \item \textbf{Dedicated brake subsystem} -- pad/rotor friction, hydraulic line compliance; currently braking has no system of its own.
    \item \textbf{Proper bearing-friction law} -- the base Coulomb/load-proportional magnitude is modeled (\S\ref{sec:bearing-friction-model}); a viscous, speed-dependent correction on top of it is not yet.
    \item \textbf{Pedal force bounds} -- what limits $\vv{P}_{3b/1,X}$ physically (max force, cadence-dependent power ceiling, left/right coupling through the crank)? none stated yet; see \S\ref{sec:pedal-force-model}.
    \item \textbf{Pedal force model selection} -- several candidate tiers are laid out (\S\ref{sec:pedal-force-model}) but which to actually use needs research/comparison against real pedal-force data before committing.
    \item \textbf{Weather-driven mass changes} -- soaked clothing, mud, water absorbed by bags/frame can meaningfully add mass; not accounted for. See \S\ref{sec:mass-inertia}.
    \item \textbf{Non-rigid rider mass properties} -- $I_k$ and CoG location are treated as constant, but a non-rigid rider (standing, leaning, shifting weight) would actually shift both; ties to the ``rider is rigidly attached'' assumption in \S\ref{sec:frame}. See \S\ref{sec:mass-inertia}.
    \item \textbf{Aerodynamic configuration argument $\xi_k$} -- needs better geometrical arguments than a single placeholder, most likely one or more angles (e.g.\ torso lean); needs research into what a rider-posture aero model actually conditions on. See \S\ref{sec:aerodynamics-model}.
    \item \textbf{Raindrop impact} -- a small, separate resistive effect distinct from drag through air; not modeled, worth investigating even though expected to be minor. See \S\ref{sec:aerodynamics-model}.
    \item \textbf{Combined-slip tire model} -- longitudinal + lateral, once turning enters 2D/3D.
    \item \textbf{Steering/trail geometry \& wheel gyroscopic effects} -- relevant once the front fork gets its own DOF, per the Timeline's 2D build order (suspension).
    \item \textbf{Road parametrization for BMX/MTB} -- does the graph form $z(x)$ used in \S\ref{sec:wheels-ground} ever actually need to fall back to a general arc-length curve $\vv{OP}(s)$? During flight there's no contact to parametrize anyway, and landing only needs $z$ at the current CoM $x$, so ordinary jumps may not require it -- true overhangs/loops (halfpipe walls) would. Unsettled; not blocking anything yet.
    \item \textbf{Contact well-posedness under high curvature} -- the single-point contact/rolling model in \S\ref{sec:wheels-ground} implicitly assumes the road's local radius of curvature stays large relative to wheel radius $r^{(2a)}$ (via $\rho(x) = (1+z'(x)^2)^{3/2}/|z''(x)|$ at the contact reference $x^*$). What tight curvature actually threatens is the \emph{geometry} -- a unique tangent point $x^*$, a well-defined frame $0$, no bridging over a dip or wedging into a crest narrower than the wheel -- not the point-force representation itself, since any pressure distribution reduces exactly to a resultant force regardless of curvature. Not yet enforced or handled as a distinct case.
\end{itemize}
\end{aiblock}


\end{document}